\chapter{Units and Measurement}
\label{ch:units}

Radio engineering spans an enormous numerical range: a weak-signal
moonbounce contact may involve powers of $10^{-21}\,\text{W}$ at the
receiver input, while a broadcast transmitter may radiate $10^{6}\,\text{W}$.
Working comfortably across eighteen or more orders of magnitude requires a
disciplined system of units and prefixes. This chapter establishes that
system before any circuit theory is introduced.

\section{The International System of Units (SI)}

Nearly all quantities used in electronics and radio are expressed in the
\textbf{International System of Units (SI)}. SI defines seven base units,
of which four are used constantly in this book:

\begin{center}
\begin{tabular}{@{}lll@{}}
\toprule
Quantity & Unit & Symbol \\
\midrule
Length & metre & m \\
Mass & kilogram & kg \\
Time & second & s \\
Electric current & ampere & A \\
\bottomrule
\end{tabular}
\end{center}

Every other electrical quantity used in this book --- voltage, resistance,
capacitance, inductance, power, frequency --- is a \emph{derived} unit,
built from these base units. For example, the volt is defined so that one
ampere flowing through one ohm of resistance produces a potential
difference of one volt; the ohm itself is defined in terms of volts and
amperes. You do not need to memorize these base-unit derivations, but
understanding that all electrical units are ultimately tied to the same
consistent system explains why formulas in this book always balance: if
you put SI units in, you get SI units out, with no hidden conversion
factors.

\begin{center}
\begin{tabular}{@{}llll@{}}
\toprule
Quantity & Unit & Symbol & In base SI units \\
\midrule
Frequency & hertz & Hz & $\text{s}^{-1}$ \\
Force & newton & N & $\text{kg}\cdot\text{m}\cdot\text{s}^{-2}$ \\
Energy & joule & J & $\text{N}\cdot\text{m}$ \\
Power & watt & W & $\text{J}/\text{s}$ \\
Charge & coulomb & C & $\text{A}\cdot\text{s}$ \\
Potential difference & volt & V & $\text{J}/\text{C}$ \\
Resistance & ohm & $\Omega$ & $\text{V}/\text{A}$ \\
Capacitance & farad & F & $\text{C}/\text{V}$ \\
Inductance & henry & H & $\text{V}\cdot\text{s}/\text{A}$ \\
Magnetic flux & weber & Wb & $\text{V}\cdot\text{s}$ \\
Magnetic flux density & tesla & T & $\text{Wb}/\text{m}^2$ \\
\bottomrule
\end{tabular}
\end{center}

\section{SI Prefixes}
\label{sec:si-prefixes}

Rather than writing very large or very small numbers with long strings of
zeros, SI defines a standard set of prefixes, each representing a power of
ten. A radio amateur must be completely fluent with the prefixes from
pico to giga; they appear on every data sheet, every dial, and in most
exam numerical questions.

\begin{center}
\begin{tabular}{@{}lllr@{}}
\toprule
Prefix & Symbol & Factor & Example \\
\midrule
tera & T & $10^{12}$ & 1\,THz (far above amateur bands) \\
giga & G & $10^{9}$ & 2.4\,GHz (Wi-Fi/amateur microwave) \\
mega & M & $10^{6}$ & 14.2\,MHz (a 20\,m band frequency) \\
kilo & k & $10^{3}$ & 3.5\,k$\Omega$ resistor \\
--- & --- & $10^{0}$ & base unit (V, A, $\Omega$, Hz, F, H, W) \\
milli & m & $10^{-3}$ & 100\,mA current \\
micro & $\mu$ & $10^{-6}$ & 100\,$\mu$F capacitor \\
nano & n & $10^{-9}$ & 10\,nH inductor \\
pico & p & $10^{-12}$ & 22\,pF trimmer capacitor \\
\bottomrule
\end{tabular}
\end{center}

\begin{keyidea}[Prefix arithmetic]
To convert between prefixes, multiply or divide by the ratio of their
powers of ten. For example, $2\,\text{mA} = 2\times10^{-3}\,\text{A} =
2000\,\mu\text{A}$. When in doubt, always convert back to the base unit
(no prefix) before doing arithmetic, then convert the answer to a
convenient prefix at the end.
\end{keyidea}

\begin{examplebox}[Converting units]
A capacitor is labeled 4700\,pF. Express this value in microfarads and in
farads.

\emph{Solution.} $1\,\text{pF} = 10^{-12}\,\text{F}$ and
$1\,\mu\text{F} = 10^{-6}\,\text{F}$, so
$1\,\text{pF} = 10^{-6}\,\mu\text{F}$. Therefore
\[
4700\,\text{pF} = 4700 \times 10^{-6}\,\mu\text{F} = 4.7\times10^{-3}\,\mu\text{F} = 0.0047\,\mu\text{F},
\]
and in base units, $4700\,\text{pF} = 4.7\times10^{-9}\,\text{F} = 4.7\,\text{nF}$.
\end{examplebox}

\section{Scientific and Engineering Notation}

\textbf{Scientific notation} writes a number as $a \times 10^{n}$ with
$1 \le a < 10$. \textbf{Engineering notation} is a variant that restricts
$n$ to multiples of 3, so the exponent always lines up with an SI prefix
--- this is the form used on virtually all electronic calculators and test
equipment. For example, $47{,}000\,\Omega$ is $4.7\times10^{4}\,\Omega$ in
scientific notation but $47\times10^{3}\,\Omega = 47\,\text{k}\Omega$ in
engineering notation. Learn to think in engineering notation; it maps
directly onto the prefix table above.

\section{Significant Figures and Measurement Uncertainty}

No physical measurement is perfectly exact. A digital multimeter reading
of ``14.235\,MHz'' implies a different precision than ``14.2\,MHz,'' even
though both could describe the same signal. The number of \textbf{significant
figures} in a stated value indicates how precisely it is known:

\begin{itemize}
\item All non-zero digits are significant.
\item Zeros between non-zero digits are significant (1002 has four).
\item Leading zeros are never significant (0.0042 has two).
\item Trailing zeros after a decimal point are significant (3.50 has three).
\end{itemize}

When multiplying or dividing measured quantities, the result should be
rounded to the same number of significant figures as the least precise
input. This matters on exams that give component values ``as measured''
rather than as nominal (ideal) values, and it matters in the workshop when
deciding how much precision a calculation actually deserves.

\section{Logarithms: The Mathematics Behind the Decibel}

Radio engineers constantly compare quantities that range over many orders
of magnitude --- transmitter power versus receiver sensitivity, for
example. The tool for compressing such ratios into manageable numbers is
the \textbf{logarithm}, formally introduced here because it underlies the
decibel system developed in Chapter~\ref{ch:decibels}.

For a base-10 logarithm,
\[
y = \log_{10}(x) \iff 10^{y} = x .
\]
Three properties make logarithms indispensable for power and gain
calculations:
\[
\log(ab) = \log a + \log b, \qquad
\log\!\left(\frac{a}{b}\right) = \log a - \log b, \qquad
\log(a^n) = n\log a .
\]
The first property is the reason decibels \emph{add} when gains and losses
are cascaded through a system (an amplifier followed by a length of lossy
cable, for instance) instead of needing to be multiplied. This single fact
is exploited throughout Chapters~\ref{ch:decibels}, \ref{ch:receivers},
and \ref{ch:transmission-lines}.

\section{Angles: Degrees and Radians}

AC circuit analysis (Chapter~\ref{ch:ac-theory} onward) uses \emph{phase
angles} extensively. Two units are used for angles: the familiar
\textbf{degree}, where a full circle is $360\degree$, and the
\textbf{radian}, where a full circle is $2\pi$ radians. Radians are the
natural unit inside trigonometric functions used in calculus-based
derivations (such as $v(t) = V_{\text{pk}}\sin(\omega t)$), while degrees
are more common in everyday circuit descriptions (``the current lags the
voltage by $30\degree$''). Conversion is direct:
\[
\theta_{\text{rad}} = \theta_{\text{deg}} \times \frac{\pi}{180}, \qquad
\theta_{\text{deg}} = \theta_{\text{rad}} \times \frac{180}{\pi}.
\]

\begin{quickreview}
\item SI base units relevant to electronics: metre, kilogram, second,
ampere; all other electrical units (volt, ohm, farad, henry, watt, \dots)
are derived from these.
\item Memorize the prefix ladder from pico ($10^{-12}$) to giga
($10^{9}$); convert to base units before doing arithmetic.
\item Engineering notation keeps exponents as multiples of 3, matching SI
prefixes directly.
\item Logarithms turn multiplication into addition --- the mathematical
basis of the decibel.
\item $360\degree = 2\pi$ radians; both units appear throughout AC theory.
\end{quickreview}

\begin{practiceproblems}
\item Convert 0.000015\,A into milliamperes and into microamperes.
\item A resistor is marked 2.2\,k$\Omega$. Express this value in ohms and
in scientific notation.
\item Convert $47\degree$ to radians, keeping three significant figures.
\item Using the properties of logarithms, simplify $\log(1000) - \log(10)$
without a calculator.
\item A frequency counter displays 7.0300\,MHz. How many significant
figures does this reading have, and what does that imply about the
counter's precision?
\end{practiceproblems}
